\scanpage{10}\mspage{\handwritten{9}}

\deleted{\unsure{L}}Ces\deleted{\unsure{e}} bicomplexes s'identifient aux complexes de \unsure{càbines|cahines|chbines|chaines} de \unsure{loggueur|longueur} 2. Donc la catégorie de $(2\text{-}\mathrm{Cat})$-groupes abéliens est équivelente à celle des complexes de chaines de longueur 2. L'explicitation de la 2-catégorie associée à un tel complexe de chaines est mon interprétation chérie. Il faudrait encore expliciter \deleted{les morphismes} 2-Grab-catégorie des homomorphismes (en un sens à préciser) entre deux\deleted{\unsure{es}} telles 2-Grab-catégories strictes, en termes d'un complexe $\mathrm{RHom}$.

L'idée qui se dégage de cette approche, et que Quillen on espère va expliciter, c'est que la \deleted{notion de catégorie de} notions de groupe \uline{abélien} dans $(n\text{-}\mathrm{Cat})$ correspond à celle de complexe de \inserted{chaines de} longueur $n$, en raisonnant soit mod. isomorphisme de part et d'autre, soit mod. une notion naturelle de $n$-équivalence d'un côté, et dans la catégorie dérivée de l'autre. Ceci doit pouvoir se préciser en utilisant des $\mathrm{Hom}^{\bullet}$ resp. des $\mathrm{RHom}^{\bullet}$.

Si maintenant $C$ est une catégorie avec un ``produit tensoriel'' \unsure{associatif} (et éventuellement unitaire), i.e. un bifoncteur avec des isomorphismes d'associativité (mais sans inverse ...), on peut lui associer \deleted{\unsure{(}}sans doute une Gr-catégorie universelle (tout comme à un monoïde on associe \unsure{un} groupe symmétrisé), d'où des invariants, qu'on pourra désigner par $K_{\unsure{o|0}}(C)$ et $K_1(C)$. On suppose qu'on pourra de même trouver, pour tout $n$, une Gr-$n$-catégorie universelle où envoyer $C$ de façon compati\unsure{bme|ble} avec $\otimes$ (\unsure{tout|dout} cela dans un sens à préciser), dont les invariants $\handwritten{\pi}_n$ seraient par définition les $K_n(C)$. \unsure{En|Un} cas intéressant est celui où $C$ est une catégorie avec produits finis (par exemple une catégorie additive) et $\otimes$ est le produit; \deleted{mais} \deleted{(}\handwritten{[}Mais on notera que dans le cas d'une catégorie non seulement \unsure{additive} mais abélienne, il y a une notion plus adéquate du $K_o$, qui devrait correspondre à une notion plus adéquate de la suite des $K_i$, inspirée du tapis de Roos; sans doute c'est en termes de (``vraie'') catégories triangulées

% archive sheet number: 9
